\section{Summary and Research Gaps}\label{sec:summary-and-research-gaps}

The systems discussed in this chapter illustrate a wide spectrum of approaches to music programming and representation,
each addressing different aspects of musical creation.
High-level composition systems such as \textbf{Sonic Pi}, \textbf{Tidal}, and \textbf{HMusic} prioritize musical
abstraction, expressiveness, and interactive composition workflows.
However, they often rely on sequential or pattern-based models that can limit explicit control over complex temporal
relationships.

Systems for real-time audio control and synthesis, including \textbf{ChucK} and \textbf{SuperCollider}, provide precise
control over timing and sound synthesis.
While these systems offer high flexibility, they operate at a relatively low level of abstraction, requiring users to
manage details related to signal processing, concurrency, and timing explicitly.

Formats for musical representation and exchange such as \textbf{MIDI} and \textbf{ABC notation} focus on encoding
musical information for storage, transmission, and playback.
While widely adopted, these formats lack support for higher-level compositional abstractions and do not provide
mechanisms for structuring music beyond sequences of events or symbolic notation.

From this analysis, several key limitations emerge.
First, many systems lack an explicit and flexible model for representing time, often relying on implicit sequencing or
pattern repetition.
Second, there is a gap between low-level audio programming languages and high-level compositional tools, with few
systems providing both expressive power and structured abstraction.
Third, existing approaches often conflate composition with execution, making it difficult to separate musical intent
from performance details.

These limitations motivate the development of the domain-specific language proposed in this thesis, which aims to
provide a timeline-based model for music composition, support deterministic and structured representations of musical
ideas, and enable modular and extensible compilation to multiple output formats.
